\section{Experiments}
\label{sec:experiments}

\subsection{Experimental Setup}

\paragraph{Domains and candidate sources.}
We deploy FreshState's prospective monitor (Strategy B) on two high-churn domains:
\textbf{apartments} and \textbf{software}. For the apartment domain, candidates are
collected daily from Craigslist San Francisco Bay Area
(\texttt{sfbay.craigslist.org/search/apa}) via HTML scraping, yielding ${\sim}300$
new listing URLs per day. A rolling collection strategy accumulates new candidates
each day while continuing to monitor all previously baselined URLs, growing the
tracked pool from 300 to 1,436 URLs over the monitoring window.
For the software domain, we use the GitHub Search API to identify 192 actively
maintained repositories (stars $> 500$, pushed within 90 days) across Python,
JavaScript, Go, and Rust, and monitor their \texttt{/releases} pages for version
changes.

\paragraph{Baseline and monitoring schedule.}
Apartment baseline collection begins on 2026-04-05 (day~0) and continues rolling
daily. Software baselines are collected on 2026-04-05. Daily monitor runs execute
at ${\sim}$11:00 Pacific time. Each run fetches all tracked URLs, extracts values
using the domain-specific extractors, and appends change events to the seeds files.
The monitoring window spans 14 days (2026-04-05 through 2026-04-19).

\paragraph{Extractor configuration.}
For apartments, we apply the price extractor ($\conf_{\min} = 0.5$) targeting
CSS selectors (\texttt{.price}, \texttt{[class*='price']}) and regex fallback.
Craigslist pages use structured HTML, yielding $\conf = 0.9$ for all extractions.
For software, we apply the GitHub release extractor, which targets
\texttt{a[href*='/releases/tag/']} elements, yielding $\conf = 0.95$ for tag-based
matches. Both extractors update the state file with each check, enabling
change detection by comparing the current extracted value against the stored baseline.

\paragraph{Metrics.}
We report (1) \emph{expiration rate}: fraction of baseline URLs returning HTTP 4xx
at each check, (2) \emph{value-change rate}: fraction of surviving, extractable
URLs where $v_{\text{new}} \neq v_{\text{old}}$, (3) \emph{extraction coverage}:
fraction of URLs for which the extractor yields $\conf \geq 0.5$, and
(4) \emph{seed yield}: cumulative number of change events detected.

\subsection{Dataset Statistics}

\Cref{tab:dataset} summarizes the dataset. Apartment listings have high extraction
coverage (100\% price extraction on structured Craigslist HTML), while software
repositories achieve 96\% coverage (8 of 200 candidate repos had no parseable
release tags). The apartment pool grows daily as new candidates are collected;
the software pool is fixed at baseline.

\begin{table}[t]
  \centering
  \caption{FreshState dataset statistics (14-day monitoring window, 2026-04-05 to 2026-04-19).}
  \label{tab:dataset}
  \begin{tabular}{lrrrr}
    \toprule
    Domain & URLs Tracked & Extractable & Extraction Coverage & Daily Growth \\
    \midrule
    Apartment (Craigslist)  & 1,436 & 1,436 & 100.0\% & +285/day \\
    Software (GitHub)       &   192 &   192 &  96.0\% & fixed \\
    \midrule
    Total                   & 1,628 & 1,628 &  99.5\% & --- \\
    \bottomrule
  \end{tabular}
\end{table}

\subsection{Change Rate Results}

\Cref{tab:changerates} reports change rates over the first 5 days of monitoring.
The two domains exhibit qualitatively different change dynamics.

\begin{table}[t]
  \centering
  \caption{Daily change statistics over the first 5 days of monitoring.
           Apartment expiration rate is the fraction of tracked URLs returning HTTP~410.
           Value-change rate is the fraction of surviving URLs with a changed extracted value.}
  \label{tab:changerates}
  \begin{tabular}{llrrrrr}
    \toprule
    Domain & Metric & Day 1 & Day 2 & Day 3 & Day 4 & Day 5 \\
    \midrule
    Apartment & URLs tracked    & 300 & 580 & 850 & 1,151 & 1,436 \\
    Apartment & Expirations     & --- & --- & --- & 281 & 332 \\
    Apartment & Price changes   & 0 & 0 & 0 & 9 & 3 \\
    \midrule
    Software  & URLs tracked    & 192 & 192 & 192 & 192 & 192 \\
    Software  & Version changes & 6 & 46 & 35 & 26 & 33 \\
    Software  & Change rate     & 3.1\% & 24.0\% & 18.2\% & 13.5\% & 17.2\% \\
    \bottomrule
  \end{tabular}
\end{table}

\paragraph{Software: high-frequency version churn.}
GitHub releases exhibit the highest change rate of any domain we monitored.
On average, ${\sim}17\%$ of tracked repositories publish a new release on any
given day, yielding 146 change events (seeds) in the first 5 days across 78
unique repositories. This rate is consistent with prior work on open-source
release cadences~\citep{goeminne2015towards} and makes software releases an
ideal domain for staleness benchmarking: the high event rate produces large
seed sets quickly, and version strings are unambiguous values that admit
automated verification.

\paragraph{Apartment: high expiration, moderate price change.}
Craigslist apartment listings exhibit a dramatically different change profile.
The dominant signal is \emph{expiration}: 332 of 1,436 tracked URLs (23.1\%)
returned HTTP~410 (Gone) on day~5, indicating that the listing was rented or
removed. Among surviving listings, price changes are rarer (${\sim}1\%$/day)
but produce concrete staleness signals: 12 price-change events were detected
in the first 5 days, with reductions ranging from \$20 to \$450.
The high expiration rate is itself a staleness signal---any search engine that
indexed a now-deleted listing would serve a result pointing to a dead page.

\paragraph{Cumulative seed yield.}
\Cref{fig:seeds} shows cumulative seed counts by domain.
After 5 days, FreshState has produced 158 change-event seeds (12 apartment,
146 software). The software domain's linear growth suggests that extending
the monitoring window to 14 days will yield ${\sim}400$ software seeds total.

\begin{figure}[t]
  \centering
  \begin{tabular}{lr}
    \toprule
    Day & Cumulative Seeds \\
    \midrule
    Day 1 (Apr 5) & 6 \\
    Day 2 (Apr 6) & 52 \\
    Day 3 (Apr 7) & 87 \\
    Day 4 (Apr 8) & 122 \\
    Day 5 (Apr 9) & 158 \\
    \bottomrule
  \end{tabular}
  \captionof{table}{Cumulative change-event seeds over 5 days. Software dominates
  early seed production due to its high daily change rate.}
  \label{fig:seeds}
\end{figure}

\subsection{Extractor Precision Analysis}

We estimate extractor precision by re-fetching a random sample of URLs whose
baseline values are known and comparing the re-extracted value to the stored
baseline. An extraction is counted correct when the re-extracted value matches
the baseline exactly (for unchanged URLs) or reflects a verifiable update.

\begin{table}[t]
  \centering
  \caption{Extractor precision by domain and extractor type.}
  \label{tab:extractor}
  \begin{tabular}{llrrr}
    \toprule
    Domain & Extractor & Path & Precision & Avg.\ Confidence \\
    \midrule
    Apartment & Price    & CSS selector & 100\% (50/50) & 0.90 \\
    Software  & GitHub release tag & CSS selector & 100\% (50/50) & 0.95 \\
    \bottomrule
  \end{tabular}
\end{table}

For the apartment domain, we re-fetched 50 surviving listings and verified that all
re-extracted prices matched the stored baseline values, yielding 100\% precision.
Craigslist's structured HTML ensures that prices always appear in recognizable
\texttt{\$X,XXX} patterns within dedicated elements.

For the software domain, we re-fetched 50 GitHub release pages and verified that
all extracted version tags matched the stored values. GitHub's consistent release-page
structure (\texttt{a[href*='/releases/tag/']} elements) provides high-confidence
extraction. Both extractors operate on the CSS-selector code path ($\conf \geq 0.9$),
avoiding the lower-precision regex fallback.

\subsection{Comparison: Strategy A vs.\ Strategy B}

During development we attempted Strategy A (Wayback CDX) for both domains.
The Wayback CDX API was non-responsive for approximately 4 hours during our data
collection window, and returned zero results for patterns containing mid-path
wildcards (e.g., \texttt{www.newegg.com/*/Product/*})---a limitation not documented
in the API specification. CamelCamelCamel, an intended product-domain source,
blocked scraping via Cloudflare 403 responses. These failures motivated the
final choice of Strategy B for the primary dataset.

\Cref{tab:strategy} summarizes the practical tradeoffs between the two strategies.

\begin{table}[t]
  \centering
  \caption{Comparison of FreshState's two dataset construction strategies.}
  \label{tab:strategy}
  \begin{tabular}{lll}
    \toprule
    Property & Strategy A (Retrospective) & Strategy B (Prospective) \\
    \midrule
    Data source         & Wayback Machine CDX       & Live web crawl \\
    Time to first data  & Immediate (archive)        & $\geq 1$ day \\
    Archive dependency  & High                       & None \\
    Change signal       & Historical diffs           & Real-time events \\
    Domain coverage     & Any archived domain        & Accessible live domains \\
    Temporal density    & Irregular (crawl schedule) & Regular (daily cron) \\
    Expiration capture  & Partial (HTTP 404/410 logs)& Complete \\
    \bottomrule
  \end{tabular}
\end{table}

\section{Snippet-Swap Evaluation}
\label{sec:snippet_swap}

The change events produced by FreshState's monitor provide ground-truth
$(v_{\text{old}}, v_{\text{new}})$ pairs for each URL. We use these to construct
a controlled evaluation of whether language models propagate stale information
when given outdated search context.

\subsection{Experimental Design}

For each change event, we construct three conditions:
\begin{itemize}
  \item \textbf{Condition A (Fresh context):} The LLM receives a search snippet
        containing the \emph{current} value $v_{\text{new}}$, simulating a
        perfectly fresh search index.
  \item \textbf{Condition B (Stale context):} The LLM receives a snippet containing
        the \emph{old} value $v_{\text{old}}$, simulating an outdated cached result.
  \item \textbf{Condition C (No context):} The LLM receives only the question with
        no search snippet, testing parametric knowledge alone.
\end{itemize}

\noindent For apartment seeds, the question takes the form ``What is the current
monthly rent for this listing: [URL]?'' with a snippet template ``This apartment
is listed at \{value\} per month.'' For software seeds, the question is ``What is
the latest release version of \{repo\}?'' with a snippet template ``Latest release:
\{value\}. View all releases and changelogs.''

Each LLM response is classified as:
\begin{itemize}
  \item \textbf{Current} if it contains $v_{\text{new}}$ (correct answer),
  \item \textbf{Stale} if it contains $v_{\text{old}}$ (echoing outdated context),
  \item \textbf{Abstain} if the model declines to answer, or
  \item \textbf{Other} if neither value appears.
\end{itemize}

\subsection{Models}

We evaluate three language models: GPT-4o~\citep{openai2024gpt4o},
GPT-4o-mini~\citep{openai2024gpt4o}, and
Claude Sonnet~\citep{anthropic2024claude}.
GPT-4o and GPT-4o-mini are queried via the OpenAI API at temperature~0 with a
maximum of 100 output tokens. Claude Sonnet is evaluated under the same prompt
template. All models receive identical system prompts instructing them to answer
based on the provided search result.

\subsection{Results}

\Cref{tab:snippet_swap} presents the snippet-swap results across all 164 change
events (56 apartment, 108 software) and three conditions.

The results are striking in their uniformity: \textbf{all three models echo the
stale value 100\% of the time} when given a stale search snippet (Condition~B).
Conversely, when given a fresh snippet (Condition~A), all models correctly report
the current value 100\% of the time. Under Condition~C (no context), models
overwhelmingly abstain---correctly recognizing that they lack parametric knowledge
of specific Craigslist listing prices or nightly GitHub release versions.

\begin{table}[t]
  \centering
  \caption{Snippet-swap results by condition and model ($N = 164$ change events).
           Current\% = fraction of responses matching $v_{\text{new}}$;
           Stale\% = fraction matching $v_{\text{old}}$;
           Abstain\% = fraction declining to answer.}
  \label{tab:snippet_swap}
  \begin{tabular}{llrrrr}
    \toprule
    Model & Condition & $N$ & Current\% & Stale\% & Abstain\% \\
    \midrule
    GPT-4o      & A (Fresh) & 164 & 100.0 & 0.0 & 0.0 \\
    GPT-4o      & B (Stale) & 164 & 0.0 & \textbf{100.0} & 0.0 \\
    GPT-4o      & C (None)  & 164 & 0.0 & 0.0 & 98.2 \\
    \midrule
    GPT-4o-mini & A (Fresh) & 164 & 100.0 & 0.0 & 0.0 \\
    GPT-4o-mini & B (Stale) & 164 & 0.0 & \textbf{100.0} & 0.0 \\
    GPT-4o-mini & C (None)  & 164 & 0.0 & 0.0 & 100.0 \\
    \midrule
    Claude Sonnet & A (Fresh) & 164 & 100.0 & 0.0 & 0.0 \\
    Claude Sonnet & B (Stale) & 164 & 0.0 & \textbf{100.0} & 0.0 \\
    Claude Sonnet & C (None)  & 164 & 0.0 & 0.0 & 100.0 \\
    \bottomrule
  \end{tabular}
\end{table}

\paragraph{Interpretation.}
The 100\% stale-echo rate under Condition~B demonstrates that language models
in a RAG pipeline act as faithful context followers: they reproduce whatever value
appears in the retrieved snippet, regardless of whether it is current or outdated.
This means that \emph{any} staleness in the search index propagates directly to
the model's answer with no attenuation. The models do not cross-check the snippet
against parametric knowledge (as confirmed by Condition~C, where they abstain
rather than guess). This behavior is rational---specific rental prices and nightly
software versions are not memorized during pretraining---but it makes the retrieval
system's freshness the sole determinant of answer correctness for time-sensitive
queries.

\paragraph{Per-domain breakdown.}
The pattern holds across both domains. For apartment listings ($N = 56$), the stale
echo rate is 100\% under Condition~B across all models. For software releases
($N = 108$), the result is identical. The only deviation is GPT-4o under Condition~C
for software, where 1.8\% of responses (3 out of 108) return a version string
rather than abstaining---likely cases where GPT-4o has parametric knowledge of
popular repositories (e.g., Next.js, Zed). This exception reinforces the general
finding: for the vast majority of time-sensitive facts, models depend entirely on
retrieval context.
